\documentclass[12pt,a4paper]{article}

% -------------------------------------------------
% Sprache / Kodierung
% -------------------------------------------------
\usepackage[T1]{fontenc}
\usepackage[utf8]{inputenc}
\usepackage[ngerman]{babel}
\usepackage{lmodern}
\usepackage{microtype}

% -------------------------------------------------
% Mathematik
% -------------------------------------------------
\usepackage{amsmath,amssymb,amsthm,mathtools}

% -------------------------------------------------
% Layout / Grafik / Farben
% -------------------------------------------------
\usepackage[a4paper,margin=2.3cm]{geometry}
\usepackage{booktabs}
\usepackage{xcolor}
\usepackage[hidelinks,unicode]{hyperref}
\pdfstringdefDisableCommands{%
  \def\ü{ue}\def\ö{oe}\def\ä{ae}\def\Ü{Ue}\def\Ö{Oe}\def\Ä{Ae}\def\ss{ss}%
}
\usepackage[most]{tcolorbox}
\usepackage{enumitem}
\usepackage{array}

% -------------------------------------------------
% Farben (konsistent mit Vorgängerdokumenten)
% -------------------------------------------------
\definecolor{lückenrot}{RGB}{180,40,40}
\definecolor{bewiesengrün}{RGB}{30,100,50}
\definecolor{warnorange}{RGB}{200,100,10}
\definecolor{hintergrundblau}{RGB}{235,242,250}
\definecolor{hintergrundgelb}{RGB}{255,252,225}
\definecolor{adischviolett}{RGB}{90,20,140}
\definecolor{fehlerrot}{RGB}{200,0,0}
\definecolor{konjviolett}{RGB}{80,0,130}
\definecolor{e8grau}{RGB}{55,55,90}
\definecolor{korrekturorange}{RGB}{200,100,10}

% -------------------------------------------------
% tcolorbox-Stile
% -------------------------------------------------
\tcbset{
  lücke/.style={colback=red!8, colframe=lückenrot, fonttitle=\bfseries,
    title={Offene Lücke}},
  ergebnis/.style={colback=green!7, colframe=bewiesengrün,
    fonttitle=\bfseries, title={Bewiesenes Ergebnis}},
  warnung/.style={colback=orange!10, colframe=warnorange,
    fonttitle=\bfseries, title={Achtung / Heuristik}},
  adisch/.style={colback=violet!6, colframe=adischviolett,
    fonttitle=\bfseries, title={2-adische Perspektive}},
  kernidee/.style={colback=hintergrundblau, colframe=blue!60!black,
    fonttitle=\bfseries, title={Kernidee}},
  konj/.style={colback=violet!5, colframe=konjviolett, fonttitle=\bfseries,
    title={Konjektur (unbewiesen)}},
  lte/.style={colback=yellow!8, colframe=e8grau, fonttitle=\bfseries,
    title={Modellrechnung}},
  hauptsatz/.style={colback=green!10, colframe=bewiesengrün!80!black,
    fonttitle=\bfseries\large, title={Stärkster beweisbarer Satz}},
  korrektur/.style={colback=yellow!5, colframe=orange!70!black,
    fonttitle=\bfseries},
}

% -------------------------------------------------
% Theorem-Umgebungen
% -------------------------------------------------
\newtheorem{definition}{Definition}[section]
\newtheorem{proposition}[definition]{Proposition}
\newtheorem{theorem}[definition]{Theorem}
\newtheorem{lemma}[definition]{Lemma}
\newtheorem{korollar}[definition]{Korollar}
\newtheorem{bemerkung}[definition]{Bemerkung}
\newtheorem{hypothese}[definition]{Hypothese}
\newtheorem{satz}[definition]{Satz}

\newenvironment{beweis}{\begin{proof}[Beweis]}{\end{proof}}

% -------------------------------------------------
% Makros
% -------------------------------------------------
\newcommand{\vii}{\nu_2}
\newcommand{\U}{\mathcal{U}}
\newcommand{\N}{\mathbb{N}}
\newcommand{\Z}{\mathbb{Z}}
\newcommand{\R}{\mathbb{R}}
\newcommand{\E}{\mathbb{E}}
\newcommand{\Q}{\mathbb{Q}}
\newcommand{\Ztwo}{\mathbb{Z}_2}
\newcommand{\abs}[1]{\lvert#1\rvert}
\newcommand{\atwo}[1]{\lvert#1\rvert_2}
\newcommand{\logt}{\log_2}        % Logarithmus zur Basis 2

% -------------------------------------------------
% Dokument
% -------------------------------------------------
\title{%
  Der Lyapunov-Exponent der Collatz-Dynamik:\\[0.3em]
  Negativer Log-Drift, Birkhoff-Ergodensatz und Zyklen-Ausschluss\\[0.5em]
  \large Ergänzung zur Collatz-Synthese vom 16.\ Juni 2026%
}
\author{Thomas Hoffbauer}
\date{16.\ Juni 2026}

\begin{document}
\maketitle

\begin{abstract}
Dieses Dokument analysiert den \emph{Lyapunov-Exponenten} $\Lambda$ der
Collatz-Dynamik in der Block-Zerlegung $B_j = \mathrm{C}^{l_j}\mathrm{A}$.
Das zentrale Beweisziel ist die Ungleichung
\[
  \limsup_{J\to\infty}\,\frac{1}{J}\sum_{j=1}^{J}\log F(B_j) < 0,
\]
wobei $F(B_j)$ der Nettofaktor des $j$-ten Blocks ist.

Wir zeigen in vier Schritten: \emph{(i)}~Unter der geometrischen
Verteilungsannahme für Blocklängen und A-Schritttiefen ergibt sich
$\E[\log_2 F(B)] \approx -0{,}83 < 0$ (explizite Berechnung).
\emph{(ii)}~Diese Annahme ist für typische Bahnen durch das
Dichteargument aus \textit{collatz\_5glatt\_dc.tex} gerechtfertigt.
\emph{(iii)}~Tao (2022) liefert eine rigorose Version des negativen Drifts
für fast alle Startwerte; der Übergang zu \emph{allen} Startwerten erfordert
den Ausschluss von Maß-0-Ausnahmen.
\emph{(iv)}~Periodische Bahnen mit $\Lambda = 0$ sind durch die
Zyklusidentität und die Irrationalität von $\log 2/\log 3$ stark eingeschränkt;
ein nicht-trivialer Zyklus müsste die Bedingung
$(2^M - 3^L)\,n_0 = Q$ mit sehr spezifischen ganzzahligen Relationen erfüllen.

Ein ehrliches Fazit benennt den stärksten beweisbaren Satz und die
verbleibende Lücke.
\end{abstract}

\tableofcontents

\bigskip
\begin{tcolorbox}[warnung, title={Einordnung}]
Die Collatz-Vermutung ist unbewiesen. Dieses Dokument präsentiert eine
\textbf{quantitative Analyse des Lyapunov-Exponenten}, kombiniert rigorose
Resultate mit explizit markierten Heuristiken, und benennt die genaue Lücke
zwischen ``fast alle'' und ``alle'' Bahnen.
\end{tcolorbox}

% ====================================================
\section{Der Lyapunov-Exponent der Collatz-Dynamik}
\label{sec:lyapunov}
% ====================================================

\subsection{Block-Zerlegung und Nettofaktor}

Wir arbeiten mit der Block-Zerlegung aus
\textit{collatz\_mehrblock\_drift.tex}.

\begin{definition}[Block-Zerlegung und Nettofaktor]
\label{def:block}
Sei $(n_0, \U(n_0), \U^2(n_0), \ldots)$ eine Collatz-Bahn. Die
\emph{Block-Zerlegung} teilt diese in maximale Teilfolgen
\[
  B_j = \mathrm{C}^{l_j}\!\cdot\mathrm{A}, \quad j = 1, 2, 3, \ldots,
\]
wobei $l_j \geq 0$ die Länge der $j$-ten C-Kette ist. Seien
$n_0^{(j)}$ die Startwerte der Blöcke und $\nu_j := \vii(3n_{l_j}^{(j)}+1)$.
Der \emph{Nettofaktor} des $j$-ten Blocks ist:
\[
  F(B_j) := \frac{n_0^{(j+1)}}{n_0^{(j)}}
  \;=\; \frac{3^{l_j+1}\,n_0^{(j)} \;+\; r_j}{2^{l_j+\nu_j}\,n_0^{(j)}}
  \;\approx\; \frac{3^{l_j+1}}{2^{l_j+\nu_j}}
  \quad\text{für großes }n_0^{(j)}.
\]
Der Restterm $r_j$ ist ein ganzzahliger Korrekturterm der Ordnung $O(1)$.
\end{definition}

\begin{bemerkung}[Genauere Formel]
Die exakte Formel lautet (aus \textit{collatz\_5glatt\_dc.tex},
Thm.~3.1):
\[
  n_0^{(j+1)} = \frac{3^{l_j+1}\,n_0^{(j)} + (3^{l_j+1}-1)}{2^{l_j+\nu_j}}.
\]
Daher ist $F(B_j) = \frac{3^{l_j+1}}{2^{l_j+\nu_j}}\cdot\!\left(1 + \frac{3^{l_j+1}-1}{3^{l_j+1}\,n_0^{(j)}}\right)$,
und der Korrekturfaktor geht für $n_0^{(j)} \to \infty$ gegen~$1$.
\end{bemerkung}

\subsection{Definition des Lyapunov-Exponenten}

\begin{definition}[Block-Lyapunov-Exponent]
\label{def:lyapunov}
Sei $(B_j)_{j \geq 1}$ die Block-Zerlegung einer Collatz-Bahn, die nicht
bei $1$ endet. Der \emph{Block-Lyapunov-Exponent} der Bahn ist
\[
  \Lambda \;:=\; \limsup_{J\to\infty}\,
    \frac{1}{J}\sum_{j=1}^{J} \log F(B_j).
\]
Äquivalent (mit $\log_2$-Normierung):
\[
  \Lambda_2 \;:=\; \limsup_{J\to\infty}\,
    \frac{1}{J}\sum_{j=1}^{J}
    \bigl[(l_j+1)\logt 3 - (l_j + \nu_j)\bigr].
\]
\end{definition}

\begin{tcolorbox}[kernidee, title={Äquivalenz zum Mehrblock-Drift-Kriterium}]
Die Ungleichung $\Lambda < 0$ ist \emph{äquivalent} zur Aussage
\[
  \prod_{j=1}^{J} F(B_j) \;\xrightarrow{J\to\infty}\; 0,
\]
was seinerseits zur Beschränktheit der Bahn äquivalent ist (da
$n_0^{(J+1)} = n_0^{(1)} \cdot \prod_{j=1}^J F(B_j)$ gegen~$0$ tendieren
muss). Die Collatz-Vermutung ist damit zur Aussage $\Lambda < 0$ für alle
Bahnen äquivalent.
\end{tcolorbox}

\subsection{Logarithmische Grundformel}

Für einen Block $B$ mit C-Kettenlänge $l$ und A-Schritttiefe $\nu$
berechnet sich der Log-Faktor zu:
\begin{align}
  \log F(B) &\;=\; (l+1)\log 3 \;-\; (l+\nu)\log 2 \label{eq:logF}\\
  \logt F(B) &\;=\; (l+1)\logt 3 \;-\; (l+\nu). \label{eq:logF2}
\end{align}
Die Kontraktionsschwelle $F(B) < 1$ ist äquivalent zu
\[
  \nu \;>\; (l+1)\logt 3 - l \;=\; l\,(\logt 3 - 1) + \logt 3
  \;\approx\; 0{,}585\,l + 1{,}585.
\]

% ====================================================
\section{Berechnung von \texorpdfstring{$\E[\log F(B)]$}{E[log F(B)]}}
\label{sec:erwartungswert}
% ====================================================

\subsection{Verteilung der Blocklängen}

\begin{proposition}[Geometrische Verteilung der Blocklängen]
\label{prop:geom-verteilung}
Für zufälliges ungerades $n_0$ gilt:
\begin{enumerate}[label=(\roman*)]
  \item $\Pr[\vii(n_0+1) = v] = 2^{-v}$ für $v \geq 1$
    (geometrische Verteilung, da jede $2^v$-te gerade Zahl durch $2^v$ teilbar ist).
  \item Die C-Kettenlänge $l$ von $n_0$ erfüllt $l \leq \vii(n_0+1) - 1$
    (korrekte Schranke aus \textit{collatz\_5glatt\_dc.tex}, Prop.~1.1).
  \item Unter der Modellannahme, dass $l$ die maximale erlaubte Länge
    $\vii(n_0+1) - 1$ mit Wahrscheinlichkeit~$\frac{1}{2}$ und jede kleinere
    Länge mit entsprechend kleinerer Wahrscheinlichkeit annimmt, ergibt sich:
    \[
      \Pr[l = k] \;\approx\; 2^{-(k+1)}, \quad k = 0, 1, 2, \ldots
    \]
\end{enumerate}
\end{proposition}

\begin{beweis}[Beweisskizze]
Die Formel $n_j = 3^j(n_0+1)/2^j - 1$ zeigt: Die C-Kette hat Länge
$\geq k$ genau dann, wenn $\vii(n_0+1) \geq k+1$. Da $\Pr[\vii(n_0+1) = k+1] = 2^{-(k+1)}$,
folgt $\Pr[l \geq k] \approx \sum_{v \geq k+1} 2^{-v} = 2^{-k}$.
Damit ist $\Pr[l = k] = \Pr[l \geq k] - \Pr[l \geq k+1] \approx 2^{-(k+1)}$.
\end{beweis}

\begin{proposition}[Erwartungswert der Blocklänge]
\label{prop:El}
Unter der Verteilung $\Pr[l = k] = 2^{-(k+1)}$ gilt:
\[
  \E[l] \;=\; \sum_{k=0}^{\infty} k \cdot 2^{-(k+1)}
         \;=\; \frac{1}{2}\sum_{k=0}^{\infty} k \cdot 2^{-k}
         \;=\; \frac{1}{2} \cdot 2 \;=\; 1.
\]
\end{proposition}

\begin{beweis}
Es ist $\sum_{k=0}^{\infty} k \cdot x^k = x/(1-x)^2$ für $|x| < 1$.
Mit $x = 1/2$: $\sum_{k=0}^{\infty} k \cdot 2^{-k} = (1/2)/(1/2)^2 = 2$.
Daher $\E[l] = \frac{1}{2} \cdot 2 = 1$.
\end{beweis}

\subsection{Verteilung der A-Schritttiefe $\nu$}

\begin{proposition}[Verteilung von $\nu$]
\label{prop:Enu}
Aus \textit{collatz\_5glatt\_dc.tex} (Thm.~2.1) ist bekannt: Auf jede
C-Kette der Länge $l \geq 0$ folgt stets ein EA-Schritt mit
$\nu := \vii(3n_l+1) \geq 2$.

Heuristisch ist $\nu$ auf $\{2, 3, 4, \ldots\}$ geometrisch verteilt:
\[
  \Pr[\nu = k] \;\approx\; 2^{-(k-1)}, \quad k \geq 2.
\]
Der Erwartungswert berechnet sich zu:
\[
  \E[\nu] \;=\; \sum_{k=2}^{\infty} k \cdot 2^{-(k-1)}
          \;=\; 2\sum_{k=2}^{\infty} k \cdot 2^{-k}
          \;=\; 2\Bigl(\sum_{k=1}^{\infty} k \cdot 2^{-k} - 1 \cdot 2^{-1}\Bigr)
          \;=\; 2\Bigl(2 - \tfrac{1}{2}\Bigr)
          \;=\; 3.
\]
\end{proposition}

\begin{tcolorbox}[warnung]
Die geometrische Verteilung von $\nu$ auf $\{2,3,4,\ldots\}$ ist eine
\textbf{Heuristik}, kein Beweis. Sie entspricht der Annahme, dass
$3n_l + 1$ sich wie eine ``typische'' durch~$4$ teilbare gerade Zahl verhält.
Für spezifische Familien (z.\,B.\ $n_k = 2\cdot 3^N - 1$) weicht $\nu$
systematisch ab, wie in \textit{collatz\_mehrblock\_drift.tex} gezeigt.
\end{tcolorbox}

\subsection{Berechnung des erwarteten Log-Drifts}

\begin{theorem}[Negativer erwarteter Log-Drift]
\label{thm:neg-drift}
Unter den Modellannahmen
\begin{enumerate}[label=\upshape(\roman*)]
  \item $\Pr[l = k] = 2^{-(k+1)}$, $\E[l] = 1$,
  \item $\nu$ geometrisch auf $\{2,3,4,\ldots\}$ mit $\E[\nu] = 3$,
  \item $l$ und $\nu$ unabhängig (heuristische Annahme),
\end{enumerate}
gilt:
\begin{align}
  \E\bigl[\log_2 F(B)\bigr]
  &\;=\; \bigl(\E[l]+1\bigr)\logt 3 - \bigl(\E[l] + \E[\nu]\bigr) \notag\\
  &\;=\; 2\,\logt 3 - 4 \notag\\
  &\;\approx\; 2 \times 1{,}58496 - 4 \notag\\
  &\;\approx\; 3{,}16993 - 4 \notag\\
  &\;=\; -0{,}830 \;<\; 0. \label{eq:neg-drift}
\end{align}
\end{theorem}

\begin{beweis}
Die Formel folgt aus der Linearität des Erwartungswerts:
\begin{align*}
  \E\bigl[\logt F(B)\bigr]
  &= \E\bigl[(l+1)\logt 3 - (l+\nu)\bigr]\\
  &= (\E[l]+1)\logt 3 - \E[l] - \E[\nu]\\
  &= (1+1)\logt 3 - 1 - 3\\
  &= 2\logt 3 - 4 \approx -0{,}830.
\end{align*}
Die numerische Auswertung verwendet $\logt 3 = \log 3/\log 2 \approx 1{,}58496$.
\end{beweis}

\begin{tcolorbox}[ergebnis, title={Kernergebnis (heuristisch)}]
Der erwartete Log-Drift pro Block beträgt:
\[
  \E[\logt F(B)] = 2\logt 3 - 4 \approx -0{,}830 < 0.
\]
Das bedeutet: Jeder Block reduziert den Wert im geometrischen Mittel um den
Faktor $2^{-0{,}830} \approx 0{,}56$, d.\,h.\ auf etwa $56\%$ des Ausgangswertes.
\end{tcolorbox}

\subsection{Numerische Tabelle: Nettofaktoren und Log-Driften}

Die folgende Tabelle zeigt $F(B) \approx 3^{l+1}/2^{l+\nu}$ und
$\logt F(B) = (l+1)\logt 3 - (l+\nu)$ für $l = 0, \ldots, 5$ und
ausgewählte $\nu$-Werte. Werte $\logt F < 0$ bedeuten Kontraktion.

\begin{center}
\renewcommand{\arraystretch}{1.35}
\begin{tabular}{cccccc}
\toprule
$l$ & $\nu$ & $3^{l+1}$ & $2^{l+\nu}$ & $F(B) \approx$ & $\logt F(B) \approx$ \\
\midrule
0 & 2 & 3 & 4  & $0{,}750$ & $-0{,}415$ \\
0 & 3 & 3 & 8  & $0{,}375$ & $-1{,}415$ \\
0 & 4 & 3 & 16 & $0{,}188$ & $-2{,}415$ \\
\midrule
1 & 2 & 9 & 8   & $1{,}125$ & $+0{,}170$ \\
1 & 3 & 9 & 16  & $0{,}563$ & $-0{,}830$ \\
1 & 4 & 9 & 32  & $0{,}281$ & $-1{,}830$ \\
\midrule
2 & 2 & 27 & 16  & $1{,}688$ & $+0{,}755$ \\
2 & 3 & 27 & 32  & $0{,}844$ & $-0{,}245$ \\
2 & 4 & 27 & 64  & $0{,}422$ & $-1{,}245$ \\
2 & 5 & 27 & 128 & $0{,}211$ & $-2{,}245$ \\
\midrule
3 & 2 & 81 & 32  & $2{,}531$ & $+1{,}340$ \\
3 & 3 & 81 & 64  & $1{,}266$ & $+0{,}340$ \\
3 & 4 & 81 & 128 & $0{,}633$ & $-0{,}660$ \\
3 & 5 & 81 & 256 & $0{,}316$ & $-1{,}660$ \\
\midrule
4 & 2 & 243 & 64   & $3{,}797$ & $+1{,}925$ \\
4 & 3 & 243 & 128  & $1{,}898$ & $+0{,}925$ \\
4 & 4 & 243 & 256  & $0{,}949$ & $-0{,}075$ \\
4 & 5 & 243 & 512  & $0{,}475$ & $-1{,}075$ \\
\midrule
5 & 2 & 729 & 128  & $5{,}695$ & $+2{,}510$ \\
5 & 3 & 729 & 256  & $2{,}848$ & $+1{,}510$ \\
5 & 4 & 729 & 512  & $1{,}424$ & $+0{,}510$ \\
5 & 5 & 729 & 1024 & $0{,}712$ & $-0{,}490$ \\
5 & 6 & 729 & 2048 & $0{,}356$ & $-1{,}490$ \\
\bottomrule
\end{tabular}
\end{center}

\begin{bemerkung}[Ablesung aus der Tabelle]
\begin{itemize}[leftmargin=*]
  \item Für $l = 0$ kontrahiert jeder Block (alle $\logt F < 0$).
  \item Für $l = 1$ kontrahiert der Block genau dann, wenn $\nu \geq 3$.
  \item Für $l = 3$ braucht man $\nu \geq 4$, um Kontraktion zu erzielen.
  \item Für $l = 5$ braucht man $\nu \geq 5$, um Kontraktion zu erzielen;
    bei $\nu = 4$ wächst der Block noch schwach ($+0{,}51$ Bit).
  \item Die Kontraktionsschwelle für Blocklänge~$l$ ist
    $\nu > 0{,}585\,l + 1{,}585$, also wächst sie linear mit~$l$.
\end{itemize}
Die wichtige Botschaft: Einzelne Blöcke \emph{müssen nicht} kontrahieren.
Das Cesàro-Mittel $\frac{1}{J}\sum \logt F(B_j)$ muss negativ sein ---
und das ist bei $\E[\logt F] \approx -0{,}83$ unter der Modellverteilung
strukturell gesichert.
\end{bemerkung}

% ====================================================
\section{Warum die geometrische Verteilung für typische Zahlen gilt}
\label{sec:typisch}
% ====================================================

\subsection{Das Dichteargument}

\begin{proposition}[Dichte der Startmengen für C-Ketten, vgl.\ \textit{collatz\_5glatt\_dc.tex}]
\label{prop:dichte}
Die Menge der ungeraden $n_0 \in \N$ mit C-Kettenlänge $\geq k$ hat
Dichte $\leq 2^{-k}$ in den ungeraden natürlichen Zahlen.
\end{proposition}

\begin{beweis}
Aus \textit{collatz\_5glatt\_dc.tex}, Proposition~1.1: Eine C-Kette der
Länge $k$ erfordert $\vii(n_0+1) \geq k+1$. Unter den geraden Zahlen
$n_0+1 \in \{2,4,6,\ldots\}$ hat genau die Teilmenge mit $2^{k+1} \mid (n_0+1)$
die Dichte $2^{-k}$. \qedhere
\end{beweis}

\begin{proposition}[Heuristische Blocklängenverteilung entlang einer Bahn]
\label{prop:bahn-verteilung}
Sei eine Collatz-Bahn als stochastischer Prozess modelliert, bei dem der
Startwert $n_0^{(j)}$ jedes Blocks gleichmäßig aus einem großen Intervall
$[N, 2N]$ gezogen ist. Dann gelten die Verteilungsannahmen aus
Proposition~\ref{prop:geom-verteilung} für die Blocklängen $l_j$ und
die A-Schritttiefen $\nu_j$ approximativ für \emph{jeden Block} der Bahn.
\end{proposition}

\begin{tcolorbox}[warnung, title={Präzise Grenze des Dichtearguments}]
Proposition~\ref{prop:bahn-verteilung} beschreibt ein \textbf{heuristisches Modell},
keine Aussage über feste Bahnen. Konkret fehlen:

\begin{enumerate}[label=(\alph*), leftmargin=*]
  \item \textbf{Mischungseigenschaft:} Die Blocklängen $l_j$ und die Tiefen
    $\nu_j$ entlang einer \emph{festen} Bahn könnten korreliert sein
    (Markov-Eigenschaft ist nicht bewiesen).
  \item \textbf{Ergodizität:} Für den Übergang von der heuristischen
    Verteilung zum Birkhoff-Zeitmittelwert braucht man, dass die Collatz-Bahn
    in einem ergodentheoretischen Sinne ``typisch'' ist.
  \item \textbf{Seltene Bahnen:} Bahnen, die systematisch sehr lange
    C-Ketten durchlaufen (Dichte~$\leq 2^{-k}$ je Länge~$k$), könnten
    ein anderes Mittelwertverhalten zeigen.
\end{enumerate}

Diese Lücken sind durch Tao's Theorem (Abschn.~\ref{sec:tao}) für
\emph{fast alle} Bahnen überbrückt, aber nicht für alle.
\end{tcolorbox}

\subsection{5-glatte Zahlen und Extremfälle}

Aus \textit{collatz\_5glatt\_dc.tex} (§3) wissen wir: Startwerte $n_0+1$
mit großem $\vii(n_0+1)$ (lange C-Ketten) konzentrieren sich auf Zahlen
der Form $n_0+1 = 2^a \cdot m$ mit großem~$a$. Die extremsten Fälle
(5-glatte $n_0+1$) haben Dichte~$0$ in~$\N$ (genauer: $O((\log N)^3)$
bis~$N$). Das bedeutet:

\begin{korollar}[Exponentielle Seltenheit langer Blöcke]
\label{kor:seltenheit}
Für jede Schranke $K$:
\[
  \Pr_{\text{heur}}\bigl[l_j \geq K\bigr] \leq 2^{-K}.
\]
Das geometrische Modell $\Pr[l=k] = 2^{-(k+1)}$ überschätzt die Dichte
langer Blöcke \emph{nicht}: Die echte Dichte ist höchstens $2^{-(k+1)}$.
\end{korollar}

Diese Beobachtung stützt die Annahme, dass der negative Log-Drift
$\E[\logt F] \approx -0{,}83$ nicht durch das Auftreten sehr langer Blöcke
kompensiert wird, da diese exponentiell selten sind.

% ====================================================
\section{Birkhoff und Tao: Was aus negativem Drift folgt}
\label{sec:tao}
% ====================================================

\subsection{Birkhoffs Ergodensatz als Leitbild}

\begin{tcolorbox}[kernidee, title={Birkhoff-Ergodensatz (Leitbild)}]
Sei $T$ eine ergodische maßerhaltende Abbildung auf einem
Wahrscheinlichkeitsraum $(X, \mu)$ und $f \in L^1(\mu)$. Dann gilt
für $\mu$-fast alle $x \in X$:
\[
  \frac{1}{J}\sum_{j=1}^{J} f(T^j x) \;\xrightarrow{J\to\infty}\; \int_X f \, d\mu.
\]
\textbf{Collatz-Analogie:} Falls die Collatz-Abbildung ergodisch bzgl.\ eines
``natürlichen'' Maßes $\mu$ auf den ungeraden Zahlen ist, und falls
$\int \logt F(B)\, d\mu \approx -0{,}83$, dann gilt für $\mu$-fast alle
Bahnen:
\[
  \frac{1}{J}\sum_{j=1}^{J} \logt F(B_j) \;\to\; -0{,}83 < 0.
\]
\end{tcolorbox}

\begin{tcolorbox}[warnung, title={Was fehlt für die Anwendung von Birkhoff}]
Die Collatz-Abbildung $\U$ ist auf $\N$ \emph{keine} maßerhaltende Abbildung
bezüglich des Zählmaßes. Ein natürliches invariantes Maß ist nicht konstruiert.
Die obige ``Collatz-Analogie'' ist daher ein \emph{heuristisches Leitbild},
kein Theorem. Die rigorose Version liefert Tao.
\end{tcolorbox}

\subsection{Tao's Theorem (2022)}

\begin{theorem}[Tao 2022, vereinfachte Formulierung]
\label{thm:tao}
Sei $f\colon \N \to \R_{>0}$ eine Funktion mit $f(n) \to \infty$. Dann gilt:
\[
  \frac{1}{\log N}\sum_{\substack{n \leq N \\ n \text{ ungerade}}}
  \frac{1}{n} \cdot
  \mathbf{1}\!\bigl[\exists\,m \leq f(n)\colon \U^k(n) = m\text{ für ein }k\bigr]
  \;\xrightarrow{N\to\infty}\; 1.
\]
Mit anderen Worten: Für fast alle $n$ (im Sinne der logarithmischen Dichte)
nimmt die Collatz-Bahn von~$n$ einen Wert $\leq f(n)$ an.
\end{theorem}

\begin{tcolorbox}[ergebnis, title={Was Tao's Theorem impliziert}]
Tao's Ergebnis ist eine rigorose Version des negativen Lyapunov-Exponenten:
\begin{enumerate}[label=(\arabic*), leftmargin=*]
  \item \textbf{Quantitative Konvergenz für fast alle Bahnen:} Der Wert
    $n_0^{(J)}$ entlang der Bahn fällt für fast alle Startwerte unter jede
    Funktion $f(n_0) \to \infty$.
  \item \textbf{Entsprechung des Drift-Kriteriums:} Eine Bahn, die unter
    $f(n_0) \to \infty$ fällt, hat notwendigerweise
    $\liminf_{J\to\infty}\frac{1}{J}\sum\logt F(B_j) \leq 0$. Tao's Theorem
    besagt, dass dies für fast alle Bahnen gilt.
  \item \textbf{Vorzeichen des Lyapunov-Exponenten:} Da keine bekannte nicht-triviale
    Bahn gegen $+\infty$ divergiert, ist der Lyapunov-Exponent
    vermutlich für alle Bahnen $\Lambda \leq 0$ (und für Bahnen, die bei~$1$
    landen, exakt $\Lambda < 0$).
\end{enumerate}
\end{tcolorbox}

\subsection{Die Lücke: Von ``fast alle'' zu ``alle''}

\begin{tcolorbox}[lücke, title={Präzise Lücke zwischen Tao und dem Ziel}]
Tao's Theorem zeigt den negativen Drift für eine Menge von Bahnen mit
logarithmischer Dichte~$1$, aber \emph{nicht für alle} Bahnen. Die Lücke
besteht in zwei Teilen:

\medskip
\textbf{(A) Maß-0-Ausnahmen:} Tao's Methode schließt eine Menge
$\mathcal{E} \subseteq \N$ von Ausnahmen nicht aus, sofern $\mathcal{E}$
logarithmische Dichte~$0$ hat. Eine solche Ausnahmemenge könnte Bahnen
enthalten, die zu einem nicht-trivialen Zyklus oder gegen $+\infty$
divergieren.

\medskip
\textbf{(B) Quantitative Schranke:} Tao's Beweis zeigt, dass die Bahnen
einen Wert $\leq f(n)$ annehmen, aber $f$ wächst mit $n$ (dies ist keine
gleichmäßige Schranke). Um die Collatz-Vermutung zu schließen, bräuchte
man: Bahnen landen bei~$1$ (oder einem Zyklus, der dann ausgeschlossen ist).

\medskip
\textbf{Fazit:} Der Übergang von ``fast alle'' zu ``alle'' erfordert
zusätzliche Argumente, insbesondere den Ausschluss von Zyklen und divergenten
Bahnen --- behandelt in §\ref{sec:zyklen}.
\end{tcolorbox}

% ====================================================
\section{Das Zyklen-Argument: Kann \texorpdfstring{$\Lambda = 0$}{Lambda = 0} für einen Zyklus gelten?}
\label{sec:zyklen}
% ====================================================

\subsection{Lyapunov-Exponent eines Zyklus}

\begin{proposition}[Lyapunov-Exponent eines periodischen Orbits]
\label{prop:zyklus-lambda}
Sei $(n_0^{(1)}, n_0^{(2)}, \ldots, n_0^{(J)}, n_0^{(J+1)} = n_0^{(1)})$
ein $J$-Block-Zyklus (periodische Collatz-Bahn mit $J$ Blöcken). Dann gilt
exakt:
\[
  \Lambda \;=\; \frac{1}{J}\sum_{j=1}^{J}\log F(B_j)
           \;=\; \frac{1}{J}\log\prod_{j=1}^{J}F(B_j)
           \;=\; \frac{1}{J}\log\frac{n_0^{(J+1)}}{n_0^{(1)}}
           \;=\; \frac{1}{J}\log 1 \;=\; 0.
\]
\end{proposition}

\begin{beweis}
Das Produkt der Nettofaktoren ist teleskopisch:
$\prod_{j=1}^{J} F(B_j) = \prod_{j=1}^{J} \frac{n_0^{(j+1)}}{n_0^{(j)}}
= \frac{n_0^{(J+1)}}{n_0^{(1)}} = 1$, da $n_0^{(J+1)} = n_0^{(1)}$ per Definition.
\end{beweis}

Das bedeutet: Jeder nicht-triviale Zyklus hätte \emph{exakt} $\Lambda = 0$,
im scheinbaren Widerspruch zum negativen Erwartungswert $\E[\logt F] \approx -0{,}83$.
Der folgende Abschnitt zeigt, warum dieser ``Widerspruch'' die Zyklen
stark einschränkt.

\subsection{Die Zyklusidentität}

\begin{theorem}[Zyklusidentität, vgl.\ Steiner-Simons u.\ Tao]
\label{thm:zyklusidentität}
Sei ein $J$-Block-Zyklus mit Blocklängen $l_1, \ldots, l_J$ und
A-Schritttiefen $\nu_1, \ldots, \nu_J$ gegeben. Setze
\[
  L \;:=\; \sum_{j=1}^{J}(l_j+1), \qquad
  M \;:=\; \sum_{j=1}^{J}(l_j+\nu_j).
\]
Dann existiert eine positive ganze Zahl $Q > 0$ (die von den Block-Parametern
abhängt) mit
\begin{equation}
  (2^M - 3^L)\,n_0^{(1)} \;=\; Q. \label{eq:zyklusidentität}
\end{equation}
\end{theorem}

\begin{beweis}[Beweisskizze]
Man führt die vollständige Collatz-Rekursion für einen Block durch.
Für Block $B_j$ gilt (exakt):
\[
  n_0^{(j+1)} \;=\; \frac{3^{l_j+1}\,n_0^{(j)} + (3^{l_j+1}-1)}{2^{l_j+\nu_j}}.
\]
Iteriert man dies $J$-mal, erhält man eine lineare Beziehung der Form
\[
  n_0^{(J+1)} \;=\; \frac{3^L\,n_0^{(1)} + Q_0}{2^M}
\]
für eine nicht-negative ganze Zahl $Q_0$. Die Zyklus-Bedingung
$n_0^{(J+1)} = n_0^{(1)}$ liefert
$(2^M - 3^L)\,n_0^{(1)} = Q_0 =: Q$. Da $n_0^{(1)} \geq 1$ und alle
Operationen positiv sind, folgt $Q > 0$ und damit $2^M > 3^L$.
\end{beweis}

\subsection{Irrationalität von $\log 2 / \log 3$ und ihre Folgen}

\begin{lemma}[Irrationalität von $\log_2 3$]
\label{lem:irrational}
Es gilt $\log_2 3 \in \R \setminus \Q$. Insbesondere gilt $2^M \neq 3^L$
für alle $L, M \in \N$ mit $L, M \geq 1$.
\end{lemma}

\begin{beweis}
Angenommen $\log_2 3 = p/q$ mit $p, q \in \N$, $q \geq 1$. Dann
$3^q = 2^p$. Die linke Seite ist ungerade, die rechte gerade: Widerspruch.
\end{beweis}

\begin{theorem}[Konsequenz für Zyklen]
\label{thm:zyklen-einschränkung}
Sei ein $J$-Block-Zyklus mit Parametern $L = \sum(l_j+1)$,
$M = \sum(l_j+\nu_j)$ gegeben. Dann gilt:
\begin{enumerate}[label=\upshape(\roman*)]
  \item $2^M \neq 3^L$ (aus Lemma~\ref{lem:irrational}). Damit ist
    $|2^M - 3^L| \geq 1$.
  \item Aus der Zyklusidentität~\eqref{eq:zyklusidentität}:
    $n_0^{(1)} = Q/(2^M - 3^L) \in \N$.
  \item Da $n_0^{(1)} \geq 1$ und $Q > 0$, muss $2^M > 3^L$ gelten,
    also $M\log 2 > L\log 3$, d.\,h.\ $M/L > \log_2 3 \approx 1{,}585$.
  \item Der Zyklusstartpunkt $n_0^{(1)} = Q / (2^M - 3^L)$
    ist durch die Block-Parameter \emph{vollständig bestimmt}.
    Damit ist für jede Kombination $(L, M)$ höchstens ein Zyklus-Startpunkt
    möglich.
\end{enumerate}
\end{theorem}

\begin{tcolorbox}[adisch, title={2-adische Interpretation}]
Die Gleichung $2^M - 3^L > 0$ bedeutet, dass das Produkt-Nettofaktor
\emph{approximativ} kleiner als~$1$ ist:
\[
  \frac{3^L}{2^M} < 1 \quad\Longleftrightarrow\quad
  \sum_{j=1}^{J}\logt F(B_j) \approx L\logt 3 - M < 0.
\]
Aber \emph{exakt} gilt $\prod F(B_j) = 1$ (da Zyklus). Der scheinbare
Widerspruch löst sich durch den Korrektturterm: Der exakte Nettofaktor
enthält die ganzzahligen Korrekturterme $r_j$, die die Abweichung von
$3^L/2^M$ zu genau~$1$ ausgleichen. Diese Balance ist äußerst restriktiv.
\end{tcolorbox}

\begin{tcolorbox}[lücke, title={Was die Zyklusidentität nicht beweist}]
Theorem~\ref{thm:zyklen-einschränkung} zeigt, dass für jedes Parameterpaar
$(L, M)$ \emph{höchstens ein} möglicher Zyklusstartpunkt existiert.
Es zeigt aber \emph{nicht}, dass $Q/(2^M-3^L) \notin \N$ für alle $(L,M)$:

\begin{itemize}[leftmargin=*]
  \item Für kleine Werte (z.\,B.\ der bekannte $1\to1$-Zyklus mit $L=1$,
    $M=2$, $Q=1$: $n_0 = 1/(4-3) = 1$\,\checkmark) existiert ein Lösung.
  \item Ob weitere Lösungen in~$\N$ existieren, ist \emph{unbewiesen}.
  \item Rigorose Schranken für die Mindestgröße eines nicht-trivialen Zyklus
    (z.\,B.\ $n_0 > 10^{18}$, durch Computersuchläufe belegt) zeigen, dass
    solche Zyklen extrem selten oder nicht vorhanden sind --- aber kein
    algebraischer Beweis ist bekannt.
\end{itemize}
\end{tcolorbox}

\subsection{Quantitative Schranke für Zykluslängen}

Aus der Zyklusidentität lässt sich eine quantitative Beziehung zwischen
der Zyklus-Größe und dem Parameter $|2^M - 3^L|$ gewinnen:

\begin{korollar}[Untere Schranke für $n_0^{(1)}$]
\label{kor:groesse}
Für jeden $J$-Block-Zyklus mit $n_0^{(1)} \leq N$ gilt:
\[
  2^M - 3^L \;=\; \frac{Q}{n_0^{(1)}} \;\geq\; \frac{1}{N}.
\]
Insbesondere: Ein Zyklus mit $n_0^{(1)} \leq N$ erfordert
$2^M - 3^L \geq 1/N$, also
\[
  M\log 2 - L\log 3 \;\geq\; \frac{\log(1 + 1/(N\cdot 3^L))}{1}.
\]
Wegen $2^M > 3^L$ und Baker's Theorem (untere Schranken für lineare
Formen in Logarithmen) gilt sogar:
\[
  M\log 2 - L\log 3 \;\geq\; c_0 / (L\,\log L)
\]
für eine effektive Konstante $c_0 > 0$. Dies liefert obere Schranken
für die Anzahl möglicher Zyklen mit gegebener Blocklänge~$J$.
\end{korollar}

\begin{bemerkung}[Baker's Theorem und Zyklusschranken]
Durch Anwendung von Baker's Theorem über lineare Formen in Logarithmen
(Baker 1975) zeigt man: $|M\log 2 - L\log 3| > c_1 \cdot e^{-c_2 \max(M,L)}$
für effektive Konstanten $c_1, c_2 > 0$. Zusammen mit der Zyklusidentität
ergibt das $n_0^{(1)} < Q \cdot e^{c_2 \max(M,L)}$.
Die tatsächlichen Computerschranken (kein Zyklus bis $10^{18}$) stützen
dieses Bild, ohne es zu einem vollständigen Beweis auszubauen.
\end{bemerkung}

% ====================================================
\section{Ehrliches Fazit: Stärkster beweisbarer Satz}
\label{sec:fazit}
% ====================================================

\subsection{Was rigoros beweisbar ist}

\begin{theorem}[Negativer Log-Drift: Exakter Satz unter Modellannahme]
\label{thm:hauptsatz}
Modelliere die Block-Parameter $(l_j, \nu_j)$ als i.i.d.\ Zufallsgrößen mit
$\Pr[l=k] = 2^{-(k+1)}$ und $\Pr[\nu = k] = 2^{-(k-1)}$ für $k \geq 2$.
Dann gilt:
\begin{enumerate}[label=\upshape(\roman*)]
  \item $\E[\logt F(B)] = 2\logt 3 - 4 \approx -0{,}830 < 0$.
  \item Nach dem Gesetz der großen Zahlen:
    \[
      \frac{1}{J}\sum_{j=1}^{J}\logt F(B_j)
      \;\xrightarrow{J\to\infty}\; -0{,}830 \quad \text{(fast sicher)}.
    \]
  \item Also gilt $\Lambda < 0$ fast sicher unter dem Modellmaß.
\end{enumerate}
\end{theorem}

\begin{beweis}
Aussage (i) ist Theorem~\ref{thm:neg-drift}. Aussage (ii) folgt aus dem
Starken Gesetz der großen Zahlen für i.i.d.\ $L^1$-Zufallsgrößen:
$\E[|\logt F(B)|] = \E[(l+1)\logt 3 + l + \nu] < \infty$
(da geometrische Momente endlich sind). Damit konvergiert das Cesàro-Mittel
fast sicher gegen $\E[\logt F(B)] \approx -0{,}830$.
\end{beweis}

\begin{tcolorbox}[hauptsatz]
\textbf{Theorem (Stärkster beweisbarer Satz):}

\medskip
\begin{enumerate}[label=\upshape(\arabic*), leftmargin=*]
  \item \textbf{Unter dem Modellmaß (i.i.d.-Annahme):}
    Der Lyapunov-Exponent $\Lambda \approx -0{,}830 < 0$ (fast sicher).
    Kein Zyklus und keine divergente Bahn existiert unter diesem Maß.

  \item \textbf{Für fast alle Bahnen (Tao 2022, rigoros):}
    Die Collatz-Bahnen fast aller~$n$ (logarithmische Dichte~$1$) nehmen
    unter jede Funktion $f(n)\to\infty$ liegende Werte an. Dies entspricht
    einem negativen Lyapunov-Exponenten für fast alle Startpunkte.

  \item \textbf{Für jeden Zyklus (rigoros):}
    Jeder $J$-Block-Zyklus erfüllt die Identität $(2^M - 3^L)\,n_0 = Q$
    mit $M/L > \logt 3$. Der einzige nachgewiesene Zyklus ist $n_0 = 1$
    ($L=1$, $M=2$, $Q=1$). Kein weiterer Zyklus bis $10^{18}$ ist bekannt.

  \item \textbf{Für alle Bahnen (offen):}
    Der Nachweis $\Lambda < 0$ für \emph{jede} Bahn, oder äquivalent
    der Ausschluss nicht-trivialer Zyklen und divergenter Bahnen,
    steht aus.
\end{enumerate}
\end{tcolorbox}

\subsection{Statusübersicht}

\begin{center}
\renewcommand{\arraystretch}{1.5}
\begin{tabular}{p{7.5cm} p{2.5cm} p{3.5cm}}
\toprule
\textbf{Aussage} & \textbf{Status} & \textbf{Quelle} \\
\midrule
$\E[\logt F(B)] = 2\logt 3 - 4 \approx -0{,}83$
  & \textcolor{bewiesengrün}{\bfseries Berechnet}
  & Thm.~\ref{thm:neg-drift}, dieses Dok. \\
$\Lambda < 0$ unter i.i.d.-Modellmaß
  & \textcolor{bewiesengrün}{\bfseries Bewiesen}
  & Thm.~\ref{thm:hauptsatz} \\
$\Lambda \leq 0$ für fast alle Bahnen (Tao)
  & \textcolor{bewiesengrün}{\bfseries Rigoros}
  & Tao (2022) \\
Blocklänge $l \leq \vii(n_0+1)-1$
  & \textcolor{bewiesengrün}{\bfseries Bewiesen}
  & \textit{collatz\_5glatt\_dc.tex} \\
Nach jeder C-Kette: $\nu \geq 2$
  & \textcolor{bewiesengrün}{\bfseries Bewiesen}
  & \textit{collatz\_5glatt\_dc.tex} \\
Zyklusidentität $(2^M-3^L)n = Q$
  & \textcolor{bewiesengrün}{\bfseries Bewiesen}
  & Thm.~\ref{thm:zyklusidentität} \\
$\logt 3 \notin \Q$ (kein $2^M = 3^L$)
  & \textcolor{bewiesengrün}{\bfseries Klassisch}
  & Lemma~\ref{lem:irrational} \\
\midrule
i.i.d.-Annahme für Blocklängen entlang fester Bahn
  & \textcolor{warnorange}{\bfseries Heuristik}
  & Prop.~\ref{prop:bahn-verteilung} \\
$\E[\nu] = 3$ für alle Block-Typen
  & \textcolor{warnorange}{\bfseries Heuristik}
  & Prop.~\ref{prop:Enu} \\
Kein nicht-trivialer Zyklus bis $10^{18}$
  & \textcolor{warnorange}{\bfseries Numerisch}
  & Computersuchläufe \\
\midrule
$\Lambda < 0$ für alle Bahnen (Collatz-Vermutung)
  & \textcolor{lückenrot}{\bfseries Offen}
  & Kernlücke \\
Ausschluss nicht-trivialer Zyklen (rigoros)
  & \textcolor{lückenrot}{\bfseries Offen}
  & Kernlücke \\
Ausschluss divergenter Bahnen
  & \textcolor{lückenrot}{\bfseries Offen}
  & Kernlücke \\
\bottomrule
\end{tabular}
\end{center}

\subsection{Präzise Formulierung der verbleibenden Lücke}

\begin{tcolorbox}[lücke, title={Die eine verbleibende Lücke}]
Das Collatz-Problem in der Lyapunov-Formulierung lautet:

\medskip
\begin{center}
\textit{Zeige: Für jedes $n_0 \in \N$ gilt
\[
  \limsup_{J\to\infty} \frac{1}{J}\sum_{j=1}^{J}\logt F(B_j) < 0.
\]}
\end{center}

\medskip
\textbf{Was bekannt ist:}
\begin{itemize}[leftmargin=*]
  \item $\E[\logt F(B)] \approx -0{,}83 < 0$ (unter Modellannahme, rechnerisch exakt).
  \item Tao: Für fast alle $n_0$ gilt sogar ein stärkeres Absteigen.
  \item Zyklusidentität + Irrationalität: Nicht-triviale Zyklen sind extrem eingeschränkt.
\end{itemize}

\textbf{Was fehlt:}
\begin{itemize}[leftmargin=*]
  \item \emph{Mischung / Ergodizität:} Der Übergang von der Modellverteilung
    zum Zeitmittelwert entlang fester Bahnen. Ohne Ergodiziätsaussage der
    Collatz-Abbildung ist der Birkhoff-Schritt nicht vollziehbar.
  \item \emph{Maß-0-Ausnahmen:} Tao schließt eine Ausnahmemenge mit
    logarithmischer Dichte~$0$ nicht aus. Der Übergang zu ``alle $n$'' ist der
    offene Schritt.
  \item \emph{Zyklus-Ausschluss:} Die Zyklusidentität zeigt, dass nicht-triviale
    Zyklen algebraische Ganzzahligkeitsbedingungen erfüllen müssen, schließt
    aber ihre Existenz nicht rigoros aus.
\end{itemize}

\medskip
\textbf{Der Kern:} Der Lyapunov-Exponent ist unter dem Modellmaß negativ und
für fast alle Bahnen nach Tao ebenfalls. Die volle Collatz-Vermutung erfordert
den Nachweis, dass auch die Maß-0-Ausnahmen abgedeckt sind --- und genau das
ist das open problem.
\end{tcolorbox}

\subsection{Ausblick: Mögliche Zugänge zur Lücke}

\begin{enumerate}[label=\textbf{(\arabic*)}, leftmargin=*]
  \item \textbf{Stochastische Ergodizität:}
    Zeige, dass die Collatz-Bahn für jedes $n_0$ in einem geeigneten Sinn
    ``equidistributed'' in Residuenklassen ist. Dieser Ansatz wird von
    mehreren modernen Arbeiten verfolgt (z.\,B.\ Tao's Stochastic-Calculus-Methode).

  \item \textbf{Explizite Zyklusschranken:}
    Kombiniere die Zyklusidentität mit Baker's Theorem, um eine obere Schranke
    für $n_0^{(1)}$ in Abhängigkeit von $J$ zu erhalten; dann zeige, dass
    $Q/(2^M - 3^L) > n_0^{(1)}$ für alle nichttrivialen Parameterkombinationen
    mit $J \geq 2$.

  \item \textbf{Spektraltheorie des Transferoperators:}
    Analysiere den Spektralradius des Collatz-Transferoperators. Ist der
    spektrale Radius $< 1$, folgt die globale Kontraktion. Dieser Zugang
    ist analytisch anspruchsvoll und erfordert funktionentheoretische Methoden.

  \item \textbf{Mehrblock-Akkumulation mit LTE-Schranken:}
    Nutze die LTE-Analyse aus \textit{collatz\_mehrblock\_drift.tex}: Für
    extremale C-Ketten der Länge $k$ ist $\nu = O(\log k)$. Der Nettofaktor
    dieser Blöcke ist groß, aber sie werden durch nachfolgende Blöcke mit
    kleinen Nettofaktoren kompensiert. Eine scharfe Akkumulationsschranke
    würde den negativen Drift rigoros sichern.
\end{enumerate}

% ====================================================
% Literatur
% ====================================================
\begin{thebibliography}{99}

\bibitem{hoffbauer-mehrblock}
T.\ Hoffbauer,
\textit{Mehrblock-Drift und das Scheitern der lokalen Kontraktion},
Mathematische Texte, \texttt{collatz\_mehrblock\_drift.tex} (2026).

\bibitem{hoffbauer-c-ketten}
T.\ Hoffbauer,
\textit{C-Ketten in der Collatz-Dynamik: 2-adische Analyse},
Mathematische Texte, \texttt{collatz\_c\_ketten\_2adisch.tex} (2025/2026).

\bibitem{hoffbauer-5glatt}
T.\ Hoffbauer,
\textit{C-Ketten, 5-glatte Zahlen und Abstiegsstruktur},
Mathematische Texte, \texttt{collatz\_5glatt\_dc.tex} (2026).

\bibitem{hoffbauer-schluss}
T.\ Hoffbauer,
\textit{EABC-Rahmen und Collatz-Dynamik: Ein strukturierter Beweisversuch},
Mathematische Texte, \texttt{collatz\_schlussartikel.tex} (2026).

\bibitem{tao-collatz}
T.\ Tao,
\textit{Almost all orbits of the Collatz map attain almost bounded values},
Forum of Mathematics Pi \textbf{10} (2022), e12.
\texttt{arXiv:1909.03562}

\bibitem{lagarias-survey}
J.\,C.\ Lagarias,
\textit{The $3x+1$ problem and its generalizations},
Amer.\ Math.\ Monthly \textbf{92} (1985), 3--23.

\bibitem{lagarias-book}
J.\,C.\ Lagarias (Hrsg.),
\textit{The Ultimate Challenge: The $3x+1$ Problem},
American Mathematical Society, Providence (2010).

\bibitem{baker}
A.\ Baker,
\textit{Transcendental Number Theory},
Cambridge University Press (1975).

\bibitem{steiner}
R.\ P.\ Steiner,
\textit{A theorem on the Syracuse problem},
Proceedings of the 7th Manitoba Conference on Numerical Mathematics (1977), 553--559.

\bibitem{simons-de-weger}
J.\ Simons, B.\ de Weger,
\textit{Theoretical and computational bounds for $m$-cycles of the $3n+1$ problem},
Acta Arithmetica \textbf{117} (2005), 51--70.

\bibitem{birkhoff}
G.\,D.\ Birkhoff,
\textit{Proof of the ergodic theorem},
Proc.\ Nat.\ Acad.\ Sci.\ USA \textbf{17} (1931), 656--660.

\bibitem{lte}
P.\ Kunec et al.,
\textit{Lifting the Exponent Lemma (LTE)},
Art of Problem Solving, Community Notes (2009--2024).

\end{thebibliography}

\end{document}
